% TOP_PROBLEM: {"id": 30007033, "title": "Hodge support on a specified K3 quadratic locus", "category_id": 5, "category": {"id": 5, "name": "algebraic_geometry", "display_name": "Algebraic Geometry", "description": "Geometric objects defined by polynomial equations.", "slug": "algebraic-geometry", "order_index": 5, "created_at": "2026-07-31T15:26:25.671Z"}, "status": "open"}
% FIRST_POSTED: 2026-10-06
% ATTEMPT_STATUS: solved
% OPENAI_CLAIM: {"schema_version":1,"source_repo":"https://github.com/openai/math","source_commit":"adc7f1241b42e322a6451854ab7e4b4c146bf78a","release_date":"2026-10-06","match":"full","resolution":"proved","independently_reviewed":false,"second_pass":"agreed","families":[{"family":"032","title":"Hodge and Kuga–Satake results for all projective K3 surfaces","subject":"Algebraic and complex geometry","manuscripts":[{"title":"Algebraic Kuga–Satake correspondences and Hodge conjectures on a K3 quadratic locus","pdf_path":"preprints/Algebraic-Kuga-Satake-correspondences-and-Hodge-conjectures-on-a-K3-quadratic-locus-September-30-2026/paper.pdf","readme_path":"preprints/Algebraic-Kuga-Satake-correspondences-and-Hodge-conjectures-on-a-K3-quadratic-locus-September-30-2026/README.md","date":"2026-09-30","theorem_ref":"main:theorem (paper.tex:179)","statement_sources":[{"path":"preprints/Algebraic-Kuga-Satake-correspondences-and-Hodge-conjectures-on-a-K3-quadratic-locus-September-30-2026/build/paper.tex","line":179,"label":"main:theorem"}]}],"lean":null,"reasoning_trace":null}]}
\documentclass[11pt]{article}
\usepackage[margin=25mm]{geometry}
\usepackage{fontspec}
\setmainfont{Latin Modern Roman}
\usepackage{amsmath,amssymb,mathtools}
\usepackage{unicode-math}
\setmathfont{Latin Modern Math}
\usepackage{enumitem}
\usepackage{xurl,hyperref,bookmark}
\hypersetup{colorlinks=true,urlcolor=blue}
\setcounter{secnumdepth}{0}
\setlength{\parindent}{0pt}
\setlength{\parskip}{5pt}
\setlength{\emergencystretch}{3em}
\begin{document}
\section*{Hodge support on a specified K3 quadratic locus}

\subsection{OpenAI's claim}
On the specified quadratic-lattice locus of projective K3 surfaces, all rational Hodge classes on every power are claimed to be algebraic; an exact algebraic Kuga--Satake tensor also gives a conditional criterion beyond that locus. The cited manuscript argues for that result; its relation to this problem is specified below.

\subsection{Mathematical statement of the claimed result}
\paragraph{Mathematical claim 1.} For a projective complex K3 surface $S$, write $T_S=\operatorname{NS}(S)_{\mathbb Q}^{\perp}$ with cup-product quadratic form $q$. Assume a rational isometric embedding
\[
(T_S,q)\hookrightarrow\mathbb U_{\mathbb Q}^{\oplus2}\perp\langle-1\rangle^4,
\]
where $\mathbb U$ is the hyperbolic plane with Gram matrix $\left(\begin{smallmatrix}0&1\\1&0\end{smallmatrix}\right)$. Choose any full even-Clifford realization $H^1(A,\mathbb Q)=C^+(T_S,q)$, any $w\in T_S$ with $q(w)\ne0$, and any rational weight-one polarization. The Kuga--Satake map is formed from $v\mapsto(a\mapsto vaw)$ and this polarization. OpenAI asserts existence of $\Gamma\in\operatorname{CH}^2(S\times A^2)_{\mathbb Q}$ inducing this exact prescribed map on $T_S$, also after transporting it to an isogenous realization. It further claims rational Hodge and generalized Hodge conjectures in all degrees of $S^m$ for every $m\ge1$. The quadratic embedding is a hypothesis, so this theorem alone is restricted to that locus.

\subsection{Outline of the argument}
Manuscript 1 (family 032, 2026-09-30): Clifford multiplication carries exterior powers of the transcendental space into degree-two cohomology of an abelian target. Transposed cycles and a polynomial inverse on the finite-dimensional image recover these summands, so divisor algebraicity supplies the required classes without an inverse Lefschetz correspondence.

\subsection{Scope relative to this problem}
Manuscript 1 (family 032, 2026-09-30): The problem selects the mathematical statement summarized here, including its stated hypotheses and parameter range. The cited theorem gives a claimed proof of that assertion in full.

\subsection{Formal verification}
Family 032 has no Lean scope page in this snapshot; no formalization of these manuscript results is listed.

This site has not executed the Lean files or checked their dependencies. These descriptions report the scope of the linked declarations, rather than a new verification of the proofs.

\subsection{Status caveat}
OpenAI attributes these results to an internal model. Its release notice cautions that results lacking full formalization can contain mistakes. This site has not independently assessed the proofs; a solved label records the scope of the reported claim.

\subsection{Sources}
\begin{itemize}
\item OpenAI, \emph{Algebraic Kuga–Satake correspondences and Hodge conjectures on a K3 quadratic locus}, 2026-09-30; citation key \texttt{OAI:Algebraic-Kuga-Satake-correspondences-and-Hodge-conjectures-on-a-K3-quadratic-locus-September-30-2026}. \href{https://github.com/openai/math/blob/adc7f1241b42e322a6451854ab7e4b4c146bf78a/preprints/Algebraic-Kuga-Satake-correspondences-and-Hodge-conjectures-on-a-K3-quadratic-locus-September-30-2026/paper.pdf}{PDF (GitHub)}; \href{https://github.com/openai/math/raw/adc7f1241b42e322a6451854ab7e4b4c146bf78a/preprints/Algebraic-Kuga-Satake-correspondences-and-Hodge-conjectures-on-a-K3-quadratic-locus-September-30-2026/paper.pdf}{PDF (direct)}; \href{https://github.com/openai/math/blob/adc7f1241b42e322a6451854ab7e4b4c146bf78a/preprints/Algebraic-Kuga-Satake-correspondences-and-Hodge-conjectures-on-a-K3-quadratic-locus-September-30-2026/README.md}{Manuscript page and citation}; \href{https://github.com/openai/math/blob/main/preprints/Algebraic-Kuga-Satake-correspondences-and-Hodge-conjectures-on-a-K3-quadratic-locus-September-30-2026/paper.pdf}{Latest version}; Statement source: \href{https://github.com/openai/math/blob/adc7f1241b42e322a6451854ab7e4b4c146bf78a/preprints/Algebraic-Kuga-Satake-correspondences-and-Hodge-conjectures-on-a-K3-quadratic-locus-September-30-2026/build/paper.tex\#L179}{main:theorem}.
\item \href{https://github.com/openai/math/blob/adc7f1241b42e322a6451854ab7e4b4c146bf78a/CONTENTS.md}{OpenAI catalogue entry}.
\end{itemize}
\end{document}
